\documentclass[10pt,a4paper]{amsart}
\usepackage[margin=19mm]{geometry}
\usepackage{amsmath,amssymb,mathtools}
\usepackage[scr=rsfs,cal=euler]{mathalfa}
\usepackage{xfrac}
\usepackage[unicode,hidelinks,bookmarks=false]{hyperref}
\newcommand{\ie}{i.e.\ }
\newcommand{\cf}{cf.\ }
\newcommand{\resp}{respectively\ }
\newcommand{\Subsection}{\subsection}
\providecommand{\nobreakdash}{\nobreak\hbox{-}}
\setlength{\emergencystretch}{2em}
\multlinegap0pt
\usepackage{amssymb}
\usepackage{bm}
\usepackage{mathtools}
\newtheorem{proposition}{Proposition}
\newcommand{\Cinf}{C^{\infty}(\M)}
\newcommand{\M}{M}
\newcommand{\XH}{\mathfrak{X}_H(\M)}
\DeclareMathOperator{\grad}{\mathrm{grad}}
\title{The Holonomy of Optimal Mass Transport:\\ The Smooth Case}
\author{Mahmoud Abdelgalil, Tryphon T. Georgiou}
\date{Source: arXiv:2608.15585v2 (CC BY 4.0)}
\begin{document}
\maketitle

\noindent\emph{Source and licence.} The Holonomy of Optimal Mass Transport:\\ The Smooth Case. Version arXiv:2608.15585v2 (CC BY 4.0), distributed by arXiv under the Creative Commons Attribution 4.0 International licence, http://creativecommons.org/licenses/by/4.0/, as recorded in the arXiv licence metadata for this exact version and shown on the abstract page https://arxiv.org/abs/2608.15585v2. Full licence notice: \texttt{licenses/arxiv-2608.15585-CC-BY-4.0.txt}.\par\smallskip

\noindent\textit{Editorial scope.} Counted unit: the article's proposition prop:bracket (source0.tex lines 231--238). The article's complete original proof of it is delivered with it (source0.tex lines 240--298, one \texttt{proof} environment of 59 lines). No mathematics is written, repaired or re-derived: the statement and the proof are the article's own lines, in the article's own notation, and no retained line is substituted or re-flowed.  No further source-local context is needed: every cross-reference of every retained line resolves inside the two delivered spans.  Nothing was omitted from any retained span, so the package carries no omission record.  No retained line contains a citation, so the wrapper carries no bibliography and no bibliography entry.  No retained line contains a figure environment, an image-inclusion command or drawing code, so no image is delivered and the package declares no asset dependency.  Editorial formatting only, in addition: an AMS \texttt{amsart} preamble that restates the article's own declarations and macros used by the retained lines, namely the environments \texttt{proposition} on separate counters, together with the standard packages those lines need.  The retained environments print in this document as Proposition 1 in order of appearance, whereas the article prints the same items with the numbering of its own full document.  The complete native source of the article as distributed in the arXiv e-print for this exact version is bundled unchanged as \texttt{sources/207763/source0.tex}.
\begin{proposition}\label{prop:bracket} Let $\Psi:\M\rightarrow\mathbb{R}^{N}$ be a smooth immersion. Then, every $X_H\in \XH$ admits the decomposition
%
$$
X_H = \sum_{i=1}^{3N} [\grad_H\phi_{2i-1}, \grad_H\phi_{2i}],
$$
%
for some collection of functions $\{\phi_i\}_{i=1}^{6N} \subset \Cinf$.
\end{proposition}

\begin{proof}
%
For any $\psi \in \Cinf$, let $\mathbb{D}(\psi)\in\mathfrak{X}(\M)$ denote the vector field defined by
%
\begin{align*}
    \mathbb{D}(\psi) := g_H(\grad_H \psi,\grad_H \psi)\, \grad_H\psi,
\end{align*}
%
and let $\{\psi_i\}_{i=1}^N$ be the component functions of the immersion $\Psi$. Define the endomorphism $P:H\rightarrow H$
%
\begin{align*}
    P:=\sum_{i=1}^{N} \mathbb{D}(\psi_i) \otimes \mathbb{D}(\psi_i)^{\flat_H},
\end{align*}
%
where $^{\flat_H}$ is the musical isomorphism induced by $g_H$. We claim that $P$ is smooth and admits a smooth inverse. To see this, note that, since $\Psi$ is an immersion, the differentials $\{\mathrm{d}\psi_i\}_{i=1}^{N}$ span the co-tangent bundle $T^{*}\M$ everywhere, and, therefore, the horizontal gradients $\{\grad_H\psi_i\}_{i=1}^{N}$ span the horizontal distribution $H$ everywhere. Recalling the definition of $\mathbb{D}(\psi)$ for a given $\psi\in\Cinf$, it is clear that the vector fields $\{\mathbb{D}(\psi_i)\}_{i=1}^{N}$ also span the horizontal distribution $H$ everywhere. Therefore, for
every $x\in\M$ and every $v\in H_x\backslash\{0\}$,
%
\begin{align*}
    g_H(x)(P(x)v,v) \;=\; \sum_{i=1}^{N} g_H(x)(\mathbb{D}(\psi_i)(x),v)^2
    \;>\; 0,
\end{align*}
%
so that $P$ is injective. Hence, $P$ is invertible and, by Cramer's rule, the inverse $P^{-1}$ is smooth. Consequently, every vector field $X_H \in \XH$ can be
decomposed as
%
\begin{align}\label{eq:wfields_decomp}
    X_H = PP^{-1}X_H=\sum_{i=1}^{N} \alpha_i\,\mathbb{D}(\psi_i),
    \qquad \alpha_i := g_H(\mathbb{D}(\psi_i),P^{-1}X_H),
\end{align}
%
where the coefficients $\alpha_i$ are smooth. We claim that each term in the summation \eqref{eq:wfields_decomp} is a linear combination of three Lie brackets of pairs of horizontal gradient fields with constant coefficients. Indeed, for all $\phi, \psi \in \Cinf$, we have that
%
\begin{equation}\label{eq:absorption}
\phi\, \mathbb{D}(\psi)
= -\tfrac{1}{4}\bigl[ \grad_H(\phi \psi^{2}),\, \grad_H \psi \bigr]
\;-\; \tfrac{1}{12}\bigl[ \grad_H \phi,\, \grad_H(\psi^{3}) \bigr]
\;-\; \tfrac{1}{4}\bigl[ \grad_H(\psi^{2}),\, \grad_H(\phi \psi) \bigr].
\end{equation}
%
To see this, note that, using the product rule for $\grad_H$ and elementary properties of the Lie Bracket, i.e.,
%
$$\grad_H(\phi \psi) = \phi\grad_H \psi + \psi \grad_H \phi, \qquad [X, \phi Y] = \phi [X, Y] + X(\phi)Y,$$ 
%
direct computation gives
%
\begin{align*}
\bigl[ \grad_H(\psi^{2}), \grad_H(\phi \psi) \bigr]
&= -2\psi^{2} [\grad_H \phi, \grad_H \psi] + 2\psi\,g_H(\grad_H \psi,\grad_H \psi) \grad_H \phi - 2\phi \,\mathbb{D}(\psi),\\
\bigl[ \grad_H \phi, \grad_H(\psi^{3}) \bigr]
&= 3\psi^{2} [\grad_H \phi, \grad_H \psi] + 6\psi\, g_H(\grad_H \phi,\grad_H \psi) \,\grad_H \psi,\\
\bigl[ \grad_H(\phi \psi^{2}), \grad_H \psi \bigr]
&= \psi^{2} [\grad_H \phi, \grad_H \psi] - 2\psi \,g_H(\grad_H \psi,\grad_H \psi)\grad_H \phi\\
&- 2\psi \, g_H(\grad_H \phi,\grad_H \psi)\grad_H \psi - 2\phi\, \mathbb{D}(\psi),
\end{align*}
%
which, after substituting into the right hand side of \eqref{eq:absorption} and simplifying, proves the identity.
%
Therefore, \eqref{eq:wfields_decomp} is a finite sum in which each term is, by virtue of \eqref{eq:absorption}, a linear combination of three Lie brackets of horizontal gradient fields with constant coefficients. Absorbing the coefficients into the functions, the conclusion follows.
\end{proof}

\end{document}
